\chapter{Analyse et spécification des besoins}
\label{chap:analyse_specification_besoins}

\section{Introduction}
\label{sec:introduction}

Dans ce chapitre, nous allons élaborer la première phase de la méthodologie Scrum, appelée également \og planification et architecture \fg{} ou \og sprint zéro \fg{}. L'objectif est d'établir une vision globale du produit, de dégager les besoins et les fonctionnalités, d'identifier les acteurs qui interagissent avec le système, d'élaborer la planification des releases, d'établir le backlog de produit et, finalement, de préparer la planification initiale des sprints et le planning de réalisation du projet.

Cette phase constitue un cadre de référence pour le projet \textit{EcoRewind}. Elle permet d'aligner les attentes des utilisateurs, les objectifs environnementaux et les choix fonctionnels avant la réalisation des incréments. La plateforme vise à encourager le tri sélectif, à faciliter l'accès à l'information locale, à valoriser les comportements écoresponsables et à outiller les acteurs territoriaux chargés de la gestion des déchets. Elle associe ainsi une dimension citoyenne, éducative, opérationnelle et territoriale.

\section{Spécification des besoins}
\label{sec:specification_besoins}

La spécification des besoins permet de déterminer les services attendus du système et les contraintes qui encadrent sa réalisation. Nous commençons par l'identification des acteurs, puis nous présentons les besoins fonctionnels associés à chaque profil ainsi que les exigences non fonctionnelles de la plateforme.

\subsection{Identification des acteurs}
\label{subsec:identification_acteurs}

L'application EcoRewind met en relation des utilisateurs ayant des responsabilités complémentaires :
\begin{itemize}
    \item \textbf{Visiteur :} personne non authentifiée qui accède aux informations publiques de la plateforme, notamment le fil d'actualité, les contenus de sensibilisation et la carte des points de collecte.
    \item \textbf{Citoyen :} utilisateur inscrit qui consulte les consignes locales, localise les points de collecte, partage ses actions, participe aux activités éducatives et cumule des récompenses grâce aux scans associés aux équipements intelligents. Il obtient également des points supplémentaires en participant aux quiz.
    \item \textbf{Éducateur :} acteur chargé de la sensibilisation environnementale. Il crée des contenus pédagogiques, des vidéos et des quiz, anime des groupes de citoyens, organise des réunions et contribue à la diffusion des bonnes pratiques de tri.
    \item \textbf{Collecteur :} acteur opérationnel chargé des activités de collecte. Il consulte les missions qui lui sont affectées, suit l'état des zones et des équipements, puis met à jour l'avancement des interventions.
    \item \textbf{Gestionnaire de point :} acteur responsable du suivi local d'un ou de plusieurs points de collecte. Il surveille leur disponibilité, leur saturation ou leur état de maintenance et participe à la remontée des informations de terrain.
    \item \textbf{Intercommunalité :} acteur institutionnel qui pilote la politique territoriale de tri. Il définit les consignes locales, centralise les données des points de collecte, organise les zones géographiques, coordonne les acteurs locaux et affecte les collecteurs aux missions.
    \item \textbf{Administrateur :} acteur chargé de la supervision globale de la plateforme. Il gère les utilisateurs, les rôles, les points de collecte, les contenus soumis par la communauté, la modération et les statistiques globales.
    \item \textbf{Super-administrateur :} administrateur disposant de privilèges renforcés pour la gestion des comptes administrateurs et la protection des rôles sensibles.
    \item \textbf{Poubelle intelligente :} équipement connecté associé à un point de collecte. Il participe à la traçabilité des opérations de tri en fournissant, selon sa configuration, des données de scan, de type de déchet, de poids ou de niveau de remplissage.
\end{itemize}

\subsection{Spécification des besoins fonctionnels par acteur}
\label{subsec:besoins_fonctionnels}

\subsubsection{Le Visiteur}

Le visiteur peut consulter le fil d'actualité, la carte des points de collecte et les contenus de sensibilisation accessibles publiquement. Il peut également créer un compte ou s'authentifier afin d'accéder aux fonctionnalités personnalisées.

\subsubsection{Le Citoyen}

Le citoyen constitue l'acteur principal de la plateforme. Il peut :
\begin{itemize}
    \item créer un compte, s'authentifier, se déconnecter et gérer les informations de son profil ;
    \item consulter les consignes locales de tri selon son territoire, sa ville ou le type de déchet ;
    \item localiser, rechercher et filtrer les points de collecte selon les déchets acceptés ou leur disponibilité ;
    \item consulter le fil communautaire, publier une action citoyenne avec texte et média, puis interagir avec les publications par des mentions d'appréciation, des commentaires ou des enregistrements ;
    \item soumettre un témoignage ou proposer un nouveau point de collecte ;
    \item recevoir les notifications relatives aux interactions, aux décisions de modération et aux activités de la plateforme ;
    \item consulter les contenus éducatifs et participer aux quiz proposés ;
    \item scanner le QR code d'une poubelle intelligente afin d'enregistrer une action de tri ;
    \item consulter son score, ses récompenses, ses badges et son impact environnemental ;
    \item échanger avec un éducateur lorsque cette communication est autorisée.
\end{itemize}

\subsubsection{L'Éducateur}

L'éducateur intervient dans la sensibilisation environnementale. Il peut créer et gérer des contenus pédagogiques, des vidéos et des quiz ; consulter les participations aux quiz ; organiser des groupes et des réunions citoyennes ; échanger avec les citoyens et participer aux actions territoriales coordonnées par l'intercommunalité.

\subsubsection{Le Collecteur}

Le collecteur peut consulter ses zones et missions, visualiser les données disponibles sur les poubelles intelligentes, mettre à jour l'état de ses affectations, renseigner des notes d'intervention et communiquer avec les acteurs habilités afin de coordonner les opérations de terrain.

\subsubsection{Le Gestionnaire de point}

Le gestionnaire de point peut visualiser les informations relatives aux points de collecte placés sous sa responsabilité, suivre leur état opérationnel -- disponible, saturé ou en maintenance --, consulter les alertes et communiquer avec les collecteurs, l'intercommunalité et les administrateurs.

\subsubsection{L'Intercommunalité}

L'intercommunalité dispose d'un espace de pilotage territorial. Elle peut :
\begin{itemize}
    \item créer, modifier, activer ou désactiver les consignes locales de tri selon le territoire, la ville et le type de déchet ;
    \item centraliser les informations relatives aux points de collecte ;
    \item consulter et coordonner les gestionnaires de point, collecteurs et éducateurs ;
    \item créer des groupes d'acteurs, leur transmettre des notifications et échanger avec eux ;
    \item créer des zones territoriales regroupant des points de collecte et des équipements ;
    \item affecter un collecteur ou un groupe de collecteurs à une zone ou à un ensemble de points ;
    \item suivre le statut, la priorité et l'avancement des missions de collecte.
\end{itemize}

\subsubsection{L'Administrateur et le Super-administrateur}

L'administrateur peut consulter le tableau de bord, gérer les comptes et les rôles selon ses autorisations, gérer les points de collecte, modérer les publications et témoignages, superviser les scans et les équipements intelligents, et consulter les statistiques de la plateforme. Le super-administrateur possède, en complément, les privilèges permettant de créer, modifier ou supprimer les comptes administrateurs et super-administrateurs.

\subsection{Spécification des besoins non fonctionnels}
\label{subsec:besoins_non_fonctionnels}

Les besoins non fonctionnels définissent les exigences de qualité du système :
\begin{itemize}
    \item \textbf{Ergonomie et accessibilité :} l'application doit proposer une navigation intuitive, cohérente et adaptée aux environnements mobiles et Web. Les écrans sont adaptés au rôle connecté afin de présenter uniquement les fonctionnalités pertinentes.
    \item \textbf{Performance :} les temps de réponse doivent rester acceptables pour la consultation du fil, de la carte et des tableaux de bord. La pagination, la compression des images et la mise en cache contribuent à cette exigence.
    \item \textbf{Sécurité :} l'accès aux ressources sensibles repose sur l'authentification, les jetons JWT et un contrôle d'accès basé sur les rôles. Les mots de passe sont hachés, les secrets sont isolés dans des variables d'environnement et l'authentification multifacteur peut être activée.
    \item \textbf{Confidentialité :} chaque utilisateur ne consulte ou ne modifie que les données auxquelles son rôle lui donne accès.
    \item \textbf{Fiabilité :} les opérations de scan, l'attribution des points et les missions de collecte doivent être traçables. Les décisions finales de modération restent sous contrôle humain.
    \item \textbf{Maintenabilité et évolutivité :} la séparation entre Flutter, routeurs API, services, schémas, repositories, modèles et migrations permet de faire évoluer la solution et d'intégrer de nouveaux territoires, rôles, déchets ou équipements.
    \item \textbf{Interopérabilité :} l'API REST assure la communication entre l'application Flutter, le backend, Firebase et le prototype matériel de poubelle intelligente.
\end{itemize}

\section{Diagramme de cas d'utilisation global}
\label{sec:diagramme_cas_utilisation_global}

Le diagramme de cas d'utilisation global présente les interactions entre les acteurs et le système EcoRewind. Il distingue les services publics accessibles au visiteur, les parcours citoyens, les fonctions éducatives, les opérations de collecte, le pilotage territorial et les fonctions d'administration.

Les services publics, tels que la consultation du fil d'actualité, de la carte des points de collecte et des contenus de sensibilisation, sont accessibles au visiteur. Les autres fonctionnalités nécessitent une authentification et sont attribuées selon le rôle de l'utilisateur. La poubelle intelligente est représentée comme un acteur externe qui transmet au système des données de scan, de type de déchet, de poids ou de niveau de remplissage.

La figure \ref{fig:diagramme_cas_utilisation_global} illustre le diagramme de cas d'utilisation global de la plateforme EcoRewind.

\begin{figure}[H]
    \centering
    \frame{\includegraphics[width=0.95\columnwidth]{image/UseCaseGlobale.png}}
    \caption{Diagramme de cas d'utilisation global d'EcoRewind}
    \label{fig:diagramme_cas_utilisation_global}
\end{figure}

\subsection{Diagramme de classes global}
\label{subsec:diagramme_classes_global}

Le diagramme de classes global présente les principales entités manipulées par le système. La classe \texttt{User} représente les utilisateurs et leurs rôles. Les classes \texttt{Post}, \texttt{Comment} et \texttt{Notification} modélisent l'espace communautaire. Les classes \texttt{CollectionPoint}, \texttt{SmartBin} et \texttt{BinScan} structurent le suivi des infrastructures de collecte et les actions de tri.

Le pilotage territorial est représenté par les classes \texttt{LocalInstruction}, \texttt{CollectorZone}, \texttt{CollectorZoneAssignment} et \texttt{CustomActorGroup}. Elles permettent de gérer les consignes locales, les zones, les affectations de collecte et la coordination des acteurs. Le diagramme met notamment en évidence qu'un citoyen peut effectuer plusieurs scans, qu'une poubelle est associée à un point de collecte et qu'un collecteur peut recevoir plusieurs affectations.

La figure \ref{fig:diagramme_classes_global} présente le diagramme de classes global de la plateforme.

\begin{figure}[H]
    \centering
    \frame{\includegraphics[width=0.95\columnwidth]{image/DiagrammeClasseGlobale.png}}
    \caption{Diagramme de classes global d'EcoRewind}
    \label{fig:diagramme_classes_global}
\end{figure}

\section{Structure et découpage de projet}
\label{sec:structure_decoupage_projet}

Cette section présente l'organisation du produit selon la démarche Scrum, le backlog priorisé, les releases, les patrons de conception et l'environnement technique du projet.

\subsection{Identification de l'équipe Scrum}
\label{subsec:equipe_scrum}

Dans un projet Scrum, l'équipe joue un rôle central dans la production continue de valeur. Les noms des membres n'étant pas définis dans le dépôt, les responsabilités attendues sont les suivantes :
\begin{itemize}
    \item \textbf{Product Owner :} porte la vision du produit, recueille les besoins des citoyens et des acteurs territoriaux, priorise le backlog et valide les incréments livrés.
    \item \textbf{Scrum Master :} accompagne l'équipe dans l'application de Scrum, facilite les cérémonies, aide à lever les obstacles et veille à l'amélioration continue.
    \item \textbf{Équipe de développement :} réalise l'application Flutter, l'API FastAPI, les modèles de données, les services Firebase, les modules de modération, les tests et les composants liés à la poubelle intelligente.
\end{itemize}

\subsection{Le Backlog du Produit global}
\label{subsec:backlog_produit_global}

Le backlog produit regroupe l'ensemble des fonctionnalités attendues de la plateforme EcoRewind. Il est organisé en épopées (\textit{Epics}) et en \textit{User Stories}, afin de faciliter la priorisation et la planification des sprints.

\setlength{\tabcolsep}{4pt}
\renewcommand{\arraystretch}{1.35}
\scriptsize
\begin{longtable}{|>{\centering\arraybackslash}p{0.8cm}|p{3.3cm}|>{\centering\arraybackslash}p{0.8cm}|p{6.7cm}|>{\centering\arraybackslash}p{1.5cm}|}
\caption{Backlog du produit EcoRewind}\label{tab:product_backlog_ecorewind}\\
\hline
\rowcolor{gray!35}
\textbf{ID} & \textbf{Fonctionnalité} \newline \textbf{(Epic)} & \textbf{ID} & \textbf{User Story} & \textbf{Priorité} \\
\hline
\endfirsthead
\multicolumn{5}{c}{\tablename\ \thetable{} -- \textit{Suite de la page précédente}} \\
\hline
\rowcolor{gray!35}
\textbf{ID} & \textbf{Fonctionnalité} \newline \textbf{(Epic)} & \textbf{ID} & \textbf{User Story} & \textbf{Priorité} \\
\hline
\endhead
\hline
\multicolumn{5}{r}{\textit{Suite à la page suivante}} \\
\endfoot
\hline
\endlastfoot

1 & Gestion des accès et des comptes utilisateurs & 1.1 & En tant que visiteur, je veux créer un compte afin d'accéder aux services personnalisés de la plateforme. & Élevée \\
\cline{3-5}
& & 1.2 & En tant qu'utilisateur, je veux m'authentifier et me déconnecter de manière sécurisée afin de protéger mon espace personnel. & Élevée \\
\cline{3-5}
& & 1.3 & En tant qu'utilisateur, je veux gérer les informations de mon profil et mon code QR afin d'être identifié par les services de tri. & Élevée \\
\cline{3-5}
& & 1.4 & En tant qu'administrateur, je veux gérer les utilisateurs et leurs rôles afin de contrôler les accès à la plateforme. & Élevée \\
\cline{3-5}
& & 1.5 & En tant que super-administrateur, je veux gérer les comptes administrateurs afin de protéger les privilèges sensibles. & Élevée \\
\hline

2 & Information citoyenne et communauté & 2.1 & En tant que visiteur, je veux consulter le fil d'actualité afin de découvrir les initiatives écologiques de la communauté. & Moyenne \\
\cline{3-5}
& & 2.2 & En tant que citoyen, je veux publier une action de tri accompagnée d'un média afin de la partager avec la communauté. & Élevée \\
\cline{3-5}
& & 2.3 & En tant que citoyen, je veux aimer, commenter et enregistrer les publications afin d'interagir avec les autres membres. & Moyenne \\
\cline{3-5}
& & 2.4 & En tant que citoyen, je veux soumettre un témoignage ou proposer un point de collecte afin de contribuer à l'amélioration du service. & Moyenne \\
\cline{3-5}
& & 2.5 & En tant qu'administrateur, je veux modérer les contenus soumis afin de garantir leur pertinence et leur sécurité. & Élevée \\
\hline

3 & Gestion des points de collecte et des consignes de tri & 3.1 & En tant que citoyen, je veux consulter les points de collecte sur une carte afin de localiser une solution adaptée à mes déchets. & Élevée \\
\cline{3-5}
& & 3.2 & En tant que citoyen, je veux filtrer les points de collecte selon le type de déchet, la localisation et la disponibilité afin de préparer mon dépôt. & Élevée \\
\cline{3-5}
& & 3.3 & En tant qu'administrateur, je veux créer, modifier et supprimer les points de collecte afin de maintenir une base d'information fiable. & Élevée \\
\cline{3-5}
& & 3.4 & En tant qu'intercommunalité, je veux gérer les consignes locales de tri selon le territoire, la ville et le type de déchet. & Élevée \\
\hline

4 & Sensibilisation environnementale & 4.1 & En tant que citoyen, je veux consulter des contenus pédagogiques afin d'améliorer mes pratiques de tri. & Moyenne \\
\cline{3-5}
& & 4.2 & En tant qu'éducateur, je veux créer et gérer des vidéos ainsi que des quiz afin de sensibiliser les citoyens. & Élevée \\
\cline{3-5}
& & 4.3 & En tant que citoyen, je veux participer à des quiz afin de renforcer mes connaissances sur le tri des déchets. & Moyenne \\
\cline{3-5}
& & 4.4 & En tant qu'éducateur, je veux organiser des groupes et des réunions afin de coordonner les activités de sensibilisation. & Moyenne \\
\hline

5 & Gestion des poubelles intelligentes et des récompenses & 5.1 & En tant que citoyen, je veux scanner le QR code d'une poubelle intelligente afin d'enregistrer mon action de tri. & Élevée \\
\cline{3-5}
& & 5.2 & En tant que système, je veux calculer automatiquement les points selon le type de déchet et, lorsque disponible, son poids. & Élevée \\
\cline{3-5}
& & 5.3 & En tant que citoyen, je veux consulter mon score, mes badges et mon impact environnemental afin de suivre mon engagement écologique. & Élevée \\
\cline{3-5}
& & 5.4 & En tant que collecteur, je veux consulter les données disponibles des poubelles intelligentes afin de mieux planifier les interventions. & Élevée \\
\cline{3-5}
& & 5.5 & En tant qu'administrateur, je veux superviser les équipements intelligents associés aux points de collecte afin d'en assurer la traçabilité. & Moyenne \\
\hline

6 & Gestion opérationnelle de la collecte & 6.1 & En tant que collecteur, je veux consulter mes zones et mes missions afin d'organiser mes tournées de collecte. & Élevée \\
\cline{3-5}
& & 6.2 & En tant que collecteur, je veux mettre à jour le statut d'une mission et ajouter une note afin de tracer son exécution. & Élevée \\
\cline{3-5}
& & 6.3 & En tant que gestionnaire de point, je veux suivre l'état des points dont j'ai la charge afin d'anticiper les besoins de collecte ou de maintenance. & Élevée \\
\hline

7 & Pilotage territorial et coordination & 7.1 & En tant qu'intercommunalité, je veux visualiser les acteurs locaux afin de coordonner les gestionnaires, les collecteurs et les éducateurs. & Élevée \\
\cline{3-5}
& & 7.2 & En tant qu'intercommunalité, je veux créer des groupes d'acteurs et leur transmettre des notifications afin de diffuser les consignes territoriales. & Moyenne \\
\cline{3-5}
& & 7.3 & En tant qu'intercommunalité, je veux créer des zones et affecter les collecteurs afin de piloter les missions de collecte sur le territoire. & Élevée \\
\cline{3-5}
& & 7.4 & En tant qu'acteur autorisé, je veux échanger par messagerie avec les profils pertinents afin de coordonner les actions. & Moyenne \\
\hline

8 & Notifications, statistiques et tableaux de bord & 8.1 & En tant qu'utilisateur, je veux recevoir des notifications liées à mes interactions, mes missions ou mes activités afin de rester informé. & Élevée \\
\cline{3-5}
& & 8.2 & En tant qu'administrateur, je veux consulter les statistiques relatives aux utilisateurs, aux scans, aux points de collecte et aux contenus afin de superviser la plateforme. & Élevée \\
\cline{3-5}
& & 8.3 & En tant qu'intercommunalité, je veux consulter les informations de pilotage territorial afin d'améliorer l'organisation de la collecte. & Élevée \\
\hline
\end{longtable}
\normalsize

\subsection{Répartition des releases}
\label{subsec:repartition_releases}

Afin de livrer progressivement de la valeur, le projet est réparti en trois releases fonctionnelles.

\begin{longtable}{|p{2.5cm}|p{11cm}|}
\hline
\rowcolor{gray!35}
\textbf{Release ID} & \textbf{Contenu de la release} \\
\hline
Release 1 : Socle citoyen & Authentification, gestion des profils et des rôles, carte des points de collecte, consignes de tri, fil communautaire, interactions et notifications initiales. \\
\hline
Release 2 : Sensibilisation, récompenses et IoT & Contenus éducatifs, quiz, groupes, témoignages, modération, poubelle intelligente, scan QR, calcul des points, badges, impact environnemental et synchronisation des données en temps réel. \\
\hline
Release 3 : Pilotage territorial et exploitation & Espaces collecteur, gestionnaire de point et intercommunalité, zones, affectations, coordination, messagerie, tableaux de bord, statistiques et stabilisation. \\
\hline
\caption{Répartition des releases du projet EcoRewind}
\label{tab:repartition_releases}
\end{longtable}

La poubelle intelligente est intégrée à la deuxième release. Ce choix permet de présenter le prototype comme une composante IoT complète, reliée au parcours citoyen, à la valorisation par les points et au suivi opérationnel de la collecte.

\subsection{Patrons de conception utilisés}
\label{subsec:patrons_conception}

\subsubsection{Architecture en couches}

Le projet sépare l'interface Flutter, les routeurs API, les services métier, les schémas, les repositories et la persistance. Cette organisation limite le couplage, améliore la lisibilité du code et facilite les tests.

\subsubsection{Data Transfer Object et validation}

Les schémas Pydantic jouent le rôle d'objets de transfert de données entre le client et l'API. Ils contrôlent la structure des requêtes et des réponses et évitent l'exposition directe des modèles de persistance.

\subsubsection{Repository}

La couche \texttt{repositories} centralise des requêtes et agrégations de données, notamment pour les statistiques. Elle isole l'accès aux données de la logique des routeurs et facilite l'évolution de la persistance.

\subsubsection{Patron État}

Les affectations de collecte suivent un cycle de vie explicite : \texttt{pending}, \texttt{in\_progress}, \texttt{done} ou \texttt{cancelled}. Les publications suivent également un cycle de modération comprenant les états de publication, d'attente de validation ou de rejet. Cette modélisation garantit la traçabilité des opérations et limite les transitions incohérentes.

\subsection{Maquette du projet}
\label{subsec:maquette_projet}

La maquette fonctionnelle d'EcoRewind se traduit par les interfaces Flutter organisées autour des principaux parcours : fil communautaire, carte des points de collecte, contenus éducatifs, récompenses, scans QR, profil, notifications, messagerie, administration et espaces métier. Les ressources visuelles et animations sont regroupées dans les répertoires \texttt{assets/images/} et \texttt{assets/lottie/}. Les captures d'écran des interfaces finales peuvent être ajoutées au rapport après leur export au format PNG ou PDF.

\subsection{Difficultés rencontrées}
\label{subsec:difficultes_rencontrees}

La réalisation d'EcoRewind réunit plusieurs domaines : développement multiplateforme, cartographie, gestion communautaire, intelligence artificielle, notification, IoT et pilotage territorial. Les principaux défis concernent la clarté des interfaces malgré la diversité des rôles, la cohérence des données entre l'API, Firebase et les équipements connectés, ainsi que la traçabilité des scans et des missions.

L'intégration de la poubelle intelligente impose de synchroniser les données issues du prototype Arduino avec les services applicatifs et de fiabiliser l'attribution des points. La modération automatique nécessite également une approche prudente : l'intelligence artificielle facilite le tri initial, tandis que la décision finale reste confiée à un administrateur.

\subsection{Environnement de travail}
\label{subsec:environnement_travail}

Le projet s'appuie sur \textbf{Flutter} et \textbf{Dart} pour l'application mobile et Web, \textbf{Python} et \textbf{FastAPI} pour l'API REST, \textbf{SQLAlchemy} et \textbf{Alembic} pour la persistance et les migrations, ainsi que \textbf{Firebase} pour les fonctions configurées de données temps réel, d'authentification et de notifications.

La cartographie utilise \textbf{OpenStreetMap} via \texttt{flutter\_map}. Le pipeline de modération exploite \textbf{PyTorch}, Text CNN, ResNet18 et NudeNet. Le prototype de poubelle intelligente s'appuie sur \textbf{Arduino} et C++. Enfin, \textbf{Git} et \textbf{GitHub} assurent le versionnement et la traçabilité des modifications.

\section{Planification des Sprints}
\label{sec:planification_sprints}

La planification suivante constitue une feuille de route indicative. Elle doit être ajustée selon la capacité réelle de l'équipe, les retours des utilisateurs et les contraintes liées au prototype matériel.

\scriptsize
\begin{longtable}{|p{2.7cm}|p{1.7cm}|p{8.1cm}|p{1.7cm}|}
\hline
\rowcolor{gray!35}
\textbf{Phase / Release} & \textbf{Sprint} & \textbf{Fonctionnalités du sprint} & \textbf{Estimation} \\
\hline
\endfirsthead
\multicolumn{4}{c}{\tablename\ \thetable{} -- \textit{Suite de la page précédente}} \\
\hline
\rowcolor{gray!35}
\textbf{Phase / Release} & \textbf{Sprint} & \textbf{Fonctionnalités du sprint} & \textbf{Estimation} \\
\hline
\endhead
\hline
\multicolumn{4}{r}{\textit{Suite à la page suivante}} \\
\endfoot
\hline
\endlastfoot
Préparation & Sprint 0 & \begin{itemize}\item Analyse des besoins et identification des acteurs.\item Élaboration du backlog produit.\item Préparation des environnements Flutter, FastAPI, base de données, Firebase et prototype Arduino.\end{itemize} & 2 semaines \\
\hline
Release 1 : Socle citoyen & Sprint 1 & \begin{itemize}\item Authentification, profils et gestion des rôles.\item Sécurité des accès et structure de navigation Flutter.\end{itemize} & 3 semaines \\
\hline
Release 1 : Socle citoyen & Sprint 2 & \begin{itemize}\item Carte des points de collecte.\item Recherche, filtrage, consignes de tri et gestion initiale des points.\end{itemize} & 3 semaines \\
\hline
Release 1 : Socle citoyen & Sprint 3 & \begin{itemize}\item Fil communautaire, téléversement de médias et interactions.\item Témoignages et notifications.\end{itemize} & 3 semaines \\
\hline
Release 2 : Sensibilisation, récompenses et IoT & Sprint 4 & \begin{itemize}\item Contenus éducatifs, vidéos et quiz.\item Groupes citoyens et réunions.\end{itemize} & 3 semaines \\
\hline
Release 2 : Sensibilisation, récompenses et IoT & Sprint 5 & \begin{itemize}\item Conception et intégration de la poubelle intelligente.\item QR code, lecture des données, communication avec Firebase ou l'API, scans et calcul des points.\end{itemize} & 3 semaines \\
\hline
Release 2 : Sensibilisation, récompenses et IoT & Sprint 6 & \begin{itemize}\item Scores, badges et impact environnemental.\item Modération automatique multicouche et validation humaine.\end{itemize} & 3 semaines \\
\hline
Release 3 : Pilotage territorial et exploitation & Sprint 7 & \begin{itemize}\item Espaces collecteur et gestionnaire de point.\item Missions, suivi d'état, alertes et notes d'intervention.\end{itemize} & 3 semaines \\
\hline
Release 3 : Pilotage territorial et exploitation & Sprint 8 & \begin{itemize}\item Espace intercommunalité : consignes locales, groupes et zones.\item Affectations, coordination et messagerie.\end{itemize} & 3 semaines \\
\hline
Release 3 : Pilotage territorial et exploitation & Sprint 9 & \begin{itemize}\item Tableaux de bord et statistiques.\item Tests fonctionnels, sécurité, optimisation et préparation du déploiement.\end{itemize} & 3 semaines \\
\hline
Clôture & -- & \begin{itemize}\item Recette finale, démonstration et documentation.\item Préparation de la livraison et rédaction du rapport de PFE.\end{itemize} & 2 semaines \\
\hline
\caption{Planification indicative des sprints du projet EcoRewind}
\label{tab:planification_sprints_ecorewind}
\end{longtable}
\normalsize

\section{Conclusion}
\label{sec:conclusion}

Ce chapitre a permis d'établir les fondations fonctionnelles et organisationnelles du projet EcoRewind. L'analyse des acteurs a mis en évidence la complémentarité entre les citoyens, les éducateurs, les collecteurs, les gestionnaires de point, l'intercommunalité et les administrateurs. Les besoins identifiés couvrent la sensibilisation environnementale, la participation citoyenne, la gestion opérationnelle de la collecte et le pilotage territorial.

Le backlog produit, les releases et la planification des sprints fournissent une feuille de route structurée pour la réalisation progressive de la plateforme. L'intégration de la poubelle intelligente dans la deuxième release consolide la dimension innovante du projet en reliant l'IoT, la valorisation des gestes de tri et l'organisation des opérations de collecte. Le chapitre suivant présentera la conception détaillée et la mise en œuvre des fonctionnalités prioritaires.
